% !PRINTED 550
\weberpage{550}
% running head: 550 Siebzehnter Abschnitt. \S.\,151.
Since all numbers from $\mathfrak{a}\mathfrak{b}$ are of the form $\Sigma\alpha\beta$, it is first of all
clear that all these numbers are divisible by $\varphi\psi$.

But if, according to \S.\,143, 4., we represent the functional $\varphi$ as the
greatest common divisor of two numbers $\alpha_1, \alpha_2$, then
these numbers, being divisible by $\varphi$, belong to the ideal $\mathfrak{a}$,
and, since a unit factor in $\varphi$ is of no
consequence, we can set

\[
\varphi = \alpha_1 x_1 + \alpha_2 x_2
\]

if $x_1, x_2$ are variables. In the same way, if
$\beta_1, \beta_2$ denote two numbers from $\mathfrak{b}$ and $y_1, y_2$ variables, we can set

\[
\psi = \beta_1 y_1 + \beta_2 y_2
\]

and from this we obtain

\[
\varphi\psi = \alpha_1 \beta_1 x_1 y_1 + \alpha_1 \beta_2 x_1 y_2 + \alpha_2 \beta_1 x_2 y_1 + \alpha_2 \beta_2 x_2 y_2.
\]

Thus by \S.\,143, 2., $\varphi\psi$ is the greatest common
divisor of the four numbers $\alpha_1\beta_1, \alpha_1\beta_2, \alpha_2\beta_1, \alpha_2\beta_2$, which belong to the ideal
$\mathfrak{a}\mathfrak{b}$, and consequently $\varphi\psi$ is the \gesperrt{greatest common
 divisor} of all numbers of the ideal $\mathfrak{a}\mathfrak{b}$.

With this the mutual assignment of the functionals and
ideals has been effected in accordance with requirements 1. to 5.

Under this assignment, \gesperrt{prime ideals} correspond to the
\gesperrt{prime functionals}, and the decomposition of the ideals into prime factors, and
in general the laws of divisibility of the ideals, result
in complete agreement with the corresponding theorems
from the theory of the functionals.

With this complete agreement it can lead to no
inconvenience if we combine the whole system of all functionals
associated with one another into one common concept
 and use for it the name \gesperrt{ideal}.

The system of all functionals associated with an integer
 is then a \gesperrt{principal ideal}.

If any number or functional is divisible by the
functionals of an ideal, we call it \gesperrt{divisible
by the ideal}; and if a functional $\varphi$ is decomposed into factors
to which the ideals $\mathfrak{a}, \mathfrak{b}, \ldots$ correspond, then,
the unit factors being omitted in the notation,
 we also set:

\weq{1}{\varphi = \mathfrak{a}\mathfrak{b} \ldots}

A basis of the functional is at the same time a basis of the ideal,
and the absolute norm of the representing functional agrees
%CONT>>
